%% Chaîne d'information : les trois blocs ACQUÉRIR / TRAITER / COMMUNIQUER.
%% Entrée : grandeurs physiques et consignes ; sortie : informations (ordres, messages).
\documentclass[tikz,border=4pt]{standalone}
\usepackage{liaisons}
\usetikzlibrary{positioning,fit}
\begin{document}
\begin{tikzpicture}[x=1cm,y=1cm,
  bloc/.style={draw=solideB, line width=1.2pt, rounded corners=4pt, fill=white,
               minimum width=2.4cm, minimum height=1.9cm, inner sep=3pt,
               align=center, font=\scriptsize},
  fleche/.style={-{Stealth[length=6pt,width=5pt]}, line width=1.8pt, draw=solideB},
  bord/.style={font=\scriptsize, text width=1.75cm, inner sep=1pt,
               execute at begin node={\hyphenpenalty=10000\exhyphenpenalty=10000}}]

%% ---- les trois blocs (même largeur imposée) --------------------------
\node[bloc] (acq) at (0,0)    {\textbf{ACQUÉRIR}\\[3pt] capteurs\\ boutons,\\ clavier};
\node[bloc] (tra) at (2.8,0) {\textbf{TRAITER}\\[3pt] carte\\ Arduino\\ automate};
\node[bloc] (com) at (5.6,0)  {\textbf{COMMUNIQUER}\\[3pt] voyants, écran\\ relais,\\ transistors\\ réseau};

%% ---- cadre englobant -------------------------------------------------
\node[draw=solideB!45, line width=0.8pt, rounded corners=7pt, dash pattern=on 3pt off 2pt,
      fit=(acq)(tra)(com), inner sep=6pt] (cadre) {};
\node[anchor=south west, font=\footnotesize\bfseries, text=solideB, inner sep=2pt]
      at (cadre.north west) {Chaîne d'information};

%% ---- liaisons internes ----------------------------------------------
\draw[fleche] (acq) -- (tra);
\draw[fleche] (tra) -- (com);

%% ---- entrée et sortie ------------------------------------------------
\draw[fleche] (-2.7,0) -- (acq.west);
\node[bord, align=right, anchor=south east] at (-1.62,0.12) {Grandeurs physiques, consignes};
\draw[fleche] (com.east) -- (8.4,0);
\node[bord, align=left, anchor=south west] at (7.22,0.12) {Informations (ordres, messages)};
\end{tikzpicture}
\end{document}
